% ==============================================================================
% Section 2: Mathematical Derivation of the Black–Scholes PDE
% Target: Page 3 (Exactly 1 Page)
% ==============================================================================

\section{Mathematical Derivation of the Black--Scholes PDE}
\label{sec:black_scholes_pde}

\subsection{Continuous Asset Dynamics and Itô's Lemma}

Let the spot price of the underlying equity $S_t$ follow a Geometric Brownian Motion (GBM) governed by the stochastic differential equation (SDE):
\begin{equation}
    \dd S_t = \mu S_t \, \dd t + \sigma S_t \, \dd W_t
    \label{eq:gbm_underlying}
\end{equation}
where $\mu$ is the expected instantaneous drift, $\sigma$ is constant volatility, and $W_t$ is a standard one-dimensional Brownian motion defined on a filtered probability space $(\Omega, \mathcal{F}, \{\mathcal{F}_t\}_{t \ge 0}, \Prob)$.

Consider a European contingent claim whose market value $V(S, t)$ is a twice continuously differentiable function of asset price $S$ and time $t$. By Itô's Lemma for two-variable functions:
\begin{equation}
    \dd V = \left( \pdiff{V}{t} + \mu S \pdiff{V}{S} + \frac{1}{2}\sigma^2 S^2 \frac{\partial^2 V}{\partial S^2} \right) \dd t + \sigma S \pdiff{V}{S} \, \dd W_t
    \label{eq:ito_claim}
\end{equation}

\subsection{Delta-Neutral Portfolio Construction and Drift Elimination}

To eliminate the stochastic noise generated by $\dd W_t$, \citet{black1973pricing} constructed a self-financing portfolio $\Pi$ consisting of one long contingent claim and a short position of $\Delta$ shares of the underlying stock:
\begin{equation}
    \Pi = V - \Delta S
    \label{eq:hedged_portfolio}
\end{equation}
Over the infinitesimal time increment $\dd t$, the change in portfolio value is:
\begin{equation}
    \dd \Pi = \dd V - \Delta \, \dd S
    \label{eq:dpi_raw}
\end{equation}
Substituting Equations~\eqref{eq:gbm_underlying} and~\eqref{eq:ito_claim} into Equation~\eqref{eq:dpi_raw}:
\begin{equation}
    \dd \Pi = \left( \pdiff{V}{t} + \mu S \pdiff{V}{S} + \frac{1}{2}\sigma^2 S^2 \frac{\partial^2 V}{\partial S^2} - \Delta \mu S \right) \dd t + \left( \sigma S \pdiff{V}{S} - \Delta \sigma S \right) \dd W_t
    \label{eq:dpi_grouped}
\end{equation}
By setting the hedge ratio equal to the instantaneous Delta:
\begin{equation}
    \Delta = \pdiff{V}{S}
    \label{eq:delta_choice}
\end{equation}
the coefficient of $\dd W_t$ vanishes identically. Crucially, the unobservable expected return drift $\mu$ also cancels out completely:
\begin{equation}
    \dd \Pi = \left( \pdiff{V}{t} + \frac{1}{2}\sigma^2 S^2 \frac{\partial^2 V}{\partial S^2} \right) \dd t
    \label{eq:dpi_deterministic}
\end{equation}

\subsection{The Black--Scholes Partial Differential Equation}

Because the hedged portfolio $\Pi$ is instantaneously riskless (its variance over $\dd t$ is zero), the principle of no-arbitrage dictates that it must earn the instantaneous risk-free rate $r$:
\begin{equation}
    \dd \Pi = r \Pi \, \dd t = r \left( V - S \pdiff{V}{S} \right) \dd t
    \label{eq:arbitrage_growth}
\end{equation}
Equating Equations~\eqref{eq:dpi_deterministic} and~\eqref{eq:arbitrage_growth} and rearranging terms yields the celebrated \textbf{Black--Scholes Partial Differential Equation}:
\begin{equation}
    \pdiff{V}{t} + r S \pdiff{V}{S} + \frac{1}{2}\sigma^2 S^2 \frac{\partial^2 V}{\partial S^2} - r V = 0
    \label{eq:bs_pde}
\end{equation}
Equation~\eqref{eq:bs_pde} is a linear parabolic PDE holding for all European derivatives, with terminal boundary conditions $V(S, T) = \max(S - K, 0)$ for European calls and $V(S, T) = \max(K - S, 0)$ for European puts.

\clearpage
